%% sections/06_ceres_algorithm.tex

\section{방향: Host-Side Slack Reclamation}
\label{sec:direction}

특성화(\S\ref{sec:results})는 두 운영 함의를 낳는다:
(A) 모델-유형 인지 방법 선택으로 회귀를 피하고,
(B) spec-decode 가 만든 host 병목을 유휴 CPU 로 떠안는다.
본 절은 (B)를 중심으로 \algname\ (host-side slack reclamation) 설계 원칙을 제시한다.
본 절은 \emph{측정된 결과가 아니라 데이터에서 도출한 방향}이며,
이를 판정할 결정적 실험(\S\ref{subsec:decisive})을 함께 정의한다.

\subsection{설계 원칙: CPU는 GPU를 흉내내지 않는다}
\label{subsec:principle}

핵심 분류 기준은 단 하나다:
\emph{CPU 가 GPU 의 일(매트멀)을 대신하면 패배, GPU 가 내려놓은 host 작업을 받으면 승리.}
전자는 H2D/D2H 동기 비용과 CPU/GPU 연산 격차(1--2 orders)로 물리적으로 진다.
후자는 spec-decode 가 step 수를 줄이며 \emph{상대적으로 커진} host-path 작업을
이미 유휴인 CPU(\S\ref{subsec:res-resource}, $\le4.9\%$)로 이전한다.

\begin{table}[t]
\centering
\caption{CPU 활용처 분류. 산 길(A/B)은 GPU 가 내려놓은 host 작업을 받고,
  죽은 길(C)은 GPU 연산을 흉내내 H2D/D2H 동기 비용에 진다.}
\label{tbl:cpu-taxonomy}
\footnotesize
\begin{tabular}{@{}lll@{}}
\toprule
 & \textbf{내용} & \textbf{근거} \\
\midrule
A1 & CPU 드래프팅(suffix tree/AMX) & GPU verify와 병렬, slack 직접 회수 \\
A2 & KV tiering (유휴 DRAM 2TB) & 배치$\uparrow$ $\to$ GPU gap 충전 \\
B1 & detokenize 오프로드 & 고-tps에서 host critical path \\
B2 & constrained-decode(문법/JSON) & code/tool 워크로드 host-heavy \\
B3 & 스케줄러/샘플링 파이프라인 & GPU feeding 지연 제거 \\
\midrule
C & GPU 매트멀 오프로드 & H2D/D2H 동기 비용에 패배(폐기) \\
\bottomrule
\end{tabular}
\end{table}

\subsection{축 A — 모델-유형 인지 방법 선택 (회귀 안전장치)}
\label{subsec:gate}

\S\ref{sec:results}는 부호의 지배 축이 \emph{모델 유형}임을 보인다.
따라서 게이트의 1차 입력은 모델 프로파일이다:
추론 특화/MoE 모델(R1 계열)은 단일 corpus 트래픽에서 vanilla,
dense 표준 모델(Llama$\cdot$Qwen$\cdot$405B-FP8)은 suffix 를 기본으로 한다.
워크로드 분류는 일부 중형 모델(Qwen-7B/32B 의 대화 대 코드)에서만 2차로 작동한다.
naive n-gram 은 실 트레이스에서 $\sim\!-90\%$ 이므로 게이트 후보에서 영구 제외한다.
이 게이트는 처리량 지배 lever 가 아니라 \emph{회귀 회피 안전장치}로 위치한다.

\subsection{축 B — host-path 회수}
\label{subsec:reclaim}

성공 셀의 GPU 여유($1-u_{\mathrm{gpu}}^{\mathrm{suf}}$)와 유휴 CPU 를 결합해
드래프팅(A1)$\cdot$detok(B1)$\cdot$문법 마스크(B2)$\cdot$스케줄(B3)을
GPU verify 와 파이프라인한다. 목표는 GPU 이용률을 다시 끌어올려
줄어든 step 수에서도 GPU 를 채우는 것이다.
이종 모델 fleet 의 multi-tenant/long-context 에서는 유휴 DRAM 2TB 를 쓰는
KV tiering(A2)으로 배치를 키워 GPU gap 을 다른 요청으로 채운다.

\subsection{결정적 실험}
\label{subsec:decisive}

축 B 의 타당성은 단일 측정으로 판정된다:
\emph{suffix ON 상태의 decode step 1개 wall-clock 을 GPU verify(커널 실행) 대
step간 gap(드래프트$\cdot$스케줄$\cdot$샘플$\cdot$detok$\cdot$동기)으로 분해}한다
(Nsight Systems inter-kernel gap + py-spy).
대상은 GPU 이용률이 가장 낮은 Qwen-7B suffix(26\%)와 대조군 Llama-70B suffix(83\%)다.
gap 이 크고 host-bound 이면(a) A1/B 가 회수 가능,
verify 내부 대역 bound 이면(b) CPU 로 불가 $\to$ A2(배치 확대)로 전환,
배치 미충전이면(c) concurrency$\uparrow$ + A2 가 1순위가 된다.
이 verdict 가 \algname\ 의 구현 우선순위를 확정한다(\S\ref{sec:discussion}).
